% bronstein_extensions.tex -- the theory behind the Bronstein completion campaign
% (src/internal/bronstein/BRONSTEIN_PLAN.md, "2026-11-14 -- Completion campaign").
%
% Every mathematical statement the implementation relies on beyond a direct
% transcription of the sources lives here: Part A derives known results in the
% exact form the code uses, Part B states and proves the extensions.  Each
% result names the routine that implements it and the regression test that
% pins it.  Build: latexmk -pdf bronstein_extensions.tex
\documentclass[11pt]{article}
\usepackage[margin=1in]{geometry}
\usepackage{amsmath,amssymb,amsthm}
\usepackage[T1]{fontenc}
\usepackage{lmodern}
\usepackage{booktabs,longtable,array}
\usepackage[hidelinks]{hyperref}

\newtheorem{theorem}{Theorem}[section]
\newtheorem{proposition}[theorem]{Proposition}
\newtheorem{lemma}[theorem]{Lemma}
\newtheorem{claim}[theorem]{Claim}
\theoremstyle{definition}
\newtheorem{definition}[theorem]{Definition}
\theoremstyle{remark}
\newtheorem{remark}[theorem]{Remark}

\newcommand{\Q}{\mathbb{Q}}
\newcommand{\Z}{\mathbb{Z}}
\newcommand{\C}{\mathbb{C}}
\newcommand{\Qbar}{\overline{\mathbb{Q}}}
\newcommand{\lc}{\operatorname{lc}}
\newcommand{\code}[1]{\texttt{#1}}

\title{Extensions of Bronstein's algorithm for one simple radical\\
\large The theory behind Mathilda's Bronstein completion campaign}
\date{Started 2026-11-14}

\begin{document}
\maketitle

\section*{Sources}
\begin{itemize}
\item[{[B]}] M.~Bronstein, \emph{Integration of elementary functions}, J.~Symbolic
  Computation 9 (1990) 117--173.
\item[{[B91]}] M.~Bronstein, \emph{The Risch differential equation on an algebraic
  curve}, ISSAC~'91.
\item[{[SI]}] M.~Bronstein, \emph{Symbolic Integration I}, 2nd ed., Springer 2005.
\item[{[T]}] B.~Trager, \emph{Integration of algebraic functions}, PhD thesis, MIT 1984.
\item[{[R]}] M.~Rosenlicht, \emph{Liouville's theorem on functions with elementary
  integrals}, Pacific J.~Math.\ 24 (1968) 153--161.
\end{itemize}

Part~A derives results the sources contain. They are written out because the
implementation depends on their exact form: the off-by-one in
\S\ref{sec:A2} was a bug before it was a derivation. Part~B states results the
sources do not contain, with proofs. Every proof in Part~B is reviewed
adversarially before its phase starts, and every claim carries the regression
test that pins it.

\paragraph{Notation.} $k$ is a differential field of characteristic~0 with
constant field $C$. $\theta$ is a monomial over $k$: either primitive,
$D\theta=\eta\in k$, or hyperexponential, $D\theta/\theta=\eta\in k$. We write
$K=k(\theta)$. For a polynomial $p$, $\lc(p)$ is its leading coefficient and
$\deg p$ its degree in $\theta$.

\part*{Part A. Derivations of known results, as implemented}
\addcontentsline{toc}{part}{Part A}

\section{Parametric logarithmic derivatives in a tower (\code{rc\_param\_log\_deriv})}
\label{sec:A1}

\paragraph{Problem.} We are given $f,\eta\in k=\Q(x)(t_0,\dots,t_{m-1})$, a
Log/Exp tower with algebraically independent monomials, and
$\theta=\exp(\int\eta)$ transcendental over $k$. Decide whether
$f=m\eta+Dz/z$ with $m\in\Z$ and $z\in k^*$. When $\eta$ is omitted, $m=0$.

\begin{claim}
Let $P$ be the irreducible factors over $\Q$ of $\operatorname{den}(f)\operatorname{den}(\eta)$,
as polynomials in $x,t_0,\dots,t_{m-1}$, excluding the exponential monomials
$t_j$. Let $E$ be the exponential monomials of $k$. Then $f=m\eta+Dz/z$ has a
solution if and only if $f$ lies in the $\Q$-span of
\[ G=\{\eta\}\cup\{Dp/p : p\in P\}\cup\{Dt_j/t_j : t_j\in E\} \]
with every coefficient an integer. The coefficients are then unique, and
$z=c\prod p^{e_p}\prod t_j^{s_j}$.
\end{claim}

\begin{proof}
Write $z=c\prod p^{e_p}\prod_{j\in E}t_j^{s_j}$, with $p$ running over the
irreducibles of $\Q[x,t]$ other than the exponential monomials. Unique
factorisation in $\Q[x,t]$ allows this, because the monomials are algebraically
independent; a logarithmic monomial is an ordinary irreducible here. Then
$Dz/z=\sum e_p\,Dp/p+\sum s_j\,Dt_j/t_j$.

Each $p$ is normal at its top variable $v$. If $v$ is $x$ or a logarithmic
monomial, then $Dv\in k_{<v}$. If $v$ is an exponential monomial other than $p$
itself, the only special irreducible is $v$. Either way $p\nmid Dp$ in
$k_{<v}[v]$.

So $Dp/p$ has a simple pole along $p$ with residue~1. No other $Dq/q$ has a
pole along $p$: those with the same top variable are coprime to $p$, those with
a lower one are constants for the partial fraction decomposition in $v$, and
$Dt_j/t_j\in k_{<t_j}$ has no pole at all.

Hence $e_p\neq0$ puts $p$ in $\operatorname{den}(m\eta+Dz/z)=\operatorname{den}(f)$,
unless the pole cancels against $m\eta$, in which case $p\mid\operatorname{den}(\eta)$.
So every $p$ with $e_p\ne0$ lies in $P$, and the representation lies in the
span of $G$ with integer coefficients.

$G$ is $\Q$-linearly independent. A relation, cleared to integers, says that
$\theta^a\prod p^b\prod t_j^c$ is a constant. For $a\neq0$ this makes $\theta$
algebraic over $k$. For $a=0$ it contradicts unique factorisation in $\Q[x,t]$.
\end{proof}

The independence of the tower is a hypothesis: a tower with $E^{x+x^2}$ beside
$E^x$ and $E^{x^2}$ breaks it. Since v0.6546 the tower builder detects this
(\code{RtTower.dependent}), and certificates on such a tower are demoted.

\section{RdeBoundDegree for a primitive monomial (\code{rde\_bound\_degree\_prim})}
\label{sec:A2}

We solve $a\,Dq+b\,q=c$ for $q\in k[\theta]$, with $D\theta=\eta\in k$. Write
$q=\sum_{i\le n}q_i\theta^i$, so that
$D(q_i\theta^i)=Dq_i\,\theta^i+i\,q_i\eta\,\theta^{i-1}$. Let $d_a,d_b,d_c$ be
the degrees of $a,b,c$.

\begin{itemize}
\item $d_b>d_a$: the top term $\lc(b)q_n\theta^{d_b+n}$ meets nothing, so
  $n\le d_c-d_b$.
\item $d_b=d_a$, with $\alpha=-\lc(b)/\lc(a)$ not of the form $Dz/z$: the
  $\theta^{d_a+n}$ coefficient $\lc(a)Dq_n+\lc(b)q_n$ vanishes if and only if
  $Dq_n/q_n=\alpha$. So $n\le d_c-d_b$.
\item $d_b=d_a$, with $\alpha=Dz/z$. Now $q_n=\gamma z$ ($\gamma$ a constant)
  cancels the top coefficient, and then $\deg(\mathrm{lhs})\le d_a+n-1$. That
  can equal $d_c$, so $n=d_c-d_a+1$ needs nothing more.

  A larger $n$ also needs the $\theta^{d_a+n-1}$ coefficient to vanish. That
  coefficient is
  \[ \lc(a)(Dq_{n-1}+n q_n\eta)+a_{d_a-1}Dq_n+\lc(b)q_{n-1}+b_{d_b-1}q_n. \]
  With $q_{n-1}=zu$ and $\lc(a)Dz+\lc(b)z=0$ it becomes
  \[ \lc(a)z\,Du+\gamma\,(n\,\lc(a)\eta z+a_{d_a-1}Dz+b_{d_b-1}z). \]
  It vanishes if and only if $-\beta=D(u/\gamma)+n\eta$, where
  $\beta=(a_{d_a-1}Dz+b_{d_b-1}z)/(z\lc(a))$. This is a limited-integration
  problem with a unique constant $s$, because two values would put $\theta$ in
  $k+C$. So
  \[ n\le\max(d_c-d_a+1,\;s), \]
  with $s$ counted only when it is a positive integer.

  The first implementation used $\max(d_c-d_a,s)$. It gave a false NoSolution
  for $y'+y/x=2\log x/x^2$, whose solution $\log^2 x/x$ has degree
  $d_c-d_a+1=2$. The regression test is in \code{tests/test\_bronstein\_rde.c}.
\item $d_b=d_a-1$: the $\theta^{d_a+n}$ coefficient $\lc(a)Dq_n$ vanishes when
  $q_n$ is a constant. The next coefficient then vanishes if and only if
  $\alpha=D(q_{n-1}/q_n)+n\eta$. So $n\le\max(d_c-d_a+1,s)$, with
  $s=\mathrm{LimitedIntegrate}(\alpha,\eta)$.
\item $d_b<d_a-1$: if the top two coefficients both vanished, $\theta$ would
  lie in $k+C$. So $n\le d_c-d_a+1$.
\end{itemize}

$\mathrm{LimitedIntegrate}(\alpha,\eta)$ against the one primitive $\theta$
asks whether $\alpha$ has an antiderivative in $K$. Any such antiderivative is
$z+s\theta+\text{const}$, and it is decided by the field integrator
(\code{rc\_limited\_integrate\_prim}). An answer that carries a new logarithm
is still reported as undecided, until ParamRischDE lands.

\section{RdeBoundDegree for a hyperexponential monomial (\code{rde\_bound\_degree\_exp})}

Here $D(q_n\theta^n)=(Dq_n+n\eta q_n)\theta^n$, so degrees are preserved.
\begin{itemize}
\item $d_a\ne d_b$: $n\le d_c-\max(d_a,d_b)$.
\item $d_a=d_b$: the top coefficient vanishes if and only if
  $-\lc(b)/\lc(a)=Dq_n/q_n+n\eta$. This is a parametric logarithmic derivative
  (\S\ref{sec:A1}) with a unique $n=m$, so $n\le\max(d_c-d_a,m)$.
\end{itemize}

\section{RdeSpecialDenomExp (\code{rde\_special\_denom\_exp})}

Suppose $q$ has a $\tau$-adic pole of order $n$. In $a\,Dq+b\,q$ its lowest
coefficient is $(a(0)(Dq_{-n}-n\eta q_{-n})+b(0)q_{-n})\tau^{-n}$. When
$\nu_\tau(b)=0$ this can vanish, if and only if $-b(0)/a(0)=Dq_{-n}/q_{-n}-n\eta$.
That is a parametric logarithmic derivative with $m=-n$. Hence
$n=\min(\nu(c)-\min(0,\nu(b)),m)$ when \S\ref{sec:A1} finds $m$.

\section{PolyRischDECancel (\code{rde\_polyrischde\_cancel})}

Here $b\in k^*$ and the equation is $Dq+bq=c$.
\begin{itemize}
\item The homogeneous equation is solvable in $K$ if and only if $-b=Dw/w$ with
  $w=z\theta^m$, which by \S\ref{sec:A1} means $-b=m\eta+Dz/z$ ($m=0$ for a
  primitive $\theta$). If it is solvable, $Dq+bq=w\,D(q/w)$, so
  $q=w\int c/w+\lambda w$. The constant $\lambda$ is fixed by the one
  coefficient $q$ must not have: a $\tau$-pole when $m<0$, or a degree beyond
  $n$.
\item If it is not solvable, each equation $Ds+(b+i\eta)s=\lc(c)$ in $k$ has at
  most one solution, since a second would give a homogeneous solution. So the
  top-down pass is forced, and a failed pass is a proof.
\item If the test is undecided, a failed pass proves nothing, and it is marked.
\end{itemize}

\section{WeakNormalizer over a tower (\code{rc\_weak\_normalizer})}

This is [SI]~\S6.1, transcribed for any monomial.
\begin{itemize}
\item Write $f=A/(E d_1)$, where $d_1$ collects the normal simple poles.
\item The residues are the roots of $R(z)=\operatorname{Res}_\theta(d_1,A-zE\,Dd_1)$.
\item Then $q=\prod_n\gcd(A-nE\,Dd_1,d_1)^n$, over the positive integer roots $n$
  of $R$.
\end{itemize}
The integer roots of $R$ are common roots of all its coefficients as a
polynomial in the tower variables. They are therefore the roots of the gcd of
those coefficients over $\Q[z]$, read off its linear factors over $\Q$.

\part*{Part B. Extensions beyond the sources}
\addcontentsline{toc}{part}{Part B}

\begin{longtable}{lp{0.5\textwidth}p{0.18\textwidth}l}
\toprule
& Result & Needed for & Status \\ \midrule
E1 & The Risch DE $z'+fz=g$ on a simple radical curve $Y^m=G$ over a
  transcendental tower $k(\theta)$: order bounds at the places over $k[\theta]$,
  at $\theta=\infty$ and at the special places; reduction of the coefficient
  system to parametric Risch DEs in $k$ & Phase 5 & open \\
E2 & Transcendental Risch over an algebraic base $k_j(y)(t_{j+1},\dots)$: the
  recursion through the differential-field record, and Liouville's theorem for a
  radical in the middle of the tower & Phase 5 & open \\
E3 & Torsion by specialisation with the special places included: goodness
  conditions, injectivity in characteristic~0 & \code{LogPart:special-genus} & v0.6548 \\
E4.1 & The constant field $\Q(i)$ with a literal $i$: every field built is a
  field over $\Q(i)$ & Phase 4.1 & v0.6551 \\
E4.1b & $i$ modulo $p\equiv1\pmod4$; Gaussian residue polynomials and $\kappa$ &
  Phase 4.1b & v0.6554 \\
E4.2 & The trigonometric front end: \code{TrigToExp} over $\Q(i)$ & Phase 4.2 & v0.6554 \\
E4.3 & Number-field constants (radicals of rationals): fields that contain them;
  certificates deferred to 4.3b & Phase 4.3 & v0.6556 \\
E4 & Symbolic parameters: the verdicts are generic, and which zero tests may be
  generic & Phase 4 & open \\
E5 & Ideals of the $d=\gcd(e,m)$ ramified places above one special point &
  Phase 3 & open \\
E6a & Reducible logarithms through a locally constant branch difference &
  \code{Tower:reducible-log} & v0.6547 \\
E6b & A cyclic radical compositum as one radical, pointwise &
  \code{Tower:cyclic-compositum} & v0.6547 \\
E6d & One radical up to a locally constant sign &
  \code{Tower:cyclic-compositum} & v0.6549 \\
E6e & The log part has no weight-$k$ coordinates, $k\ne1$ (gap \#12) &
  \code{LogPart:weights} & v0.6550 \\
E6c & Exponents with a rational log-linear part &
  \code{Tower:exp-of-log-radical} & v0.6547 \\
\bottomrule
\end{longtable}

\begin{remark}[{[B] pp.~155--157}]
This case needs $\eta'=\sum e_i\gamma_i'/\gamma_i$ with constants $e_i$. With
rational $e_i$, $\theta$ would be algebraic over $k$, so over the constant
field $\Q$ the case cannot occur in a valid tower. It does occur for irrational
or transcendental exponents ($x^{\sqrt2}$, $x^a$), so it belongs to Phase~4.
The same reading refuted a suspected false certificate in
\code{brIntegrateInKLogLike}: Def.~2.21 of [B] allows logarithms of any element
of $k$.
\end{remark}

\section*{E3. Torsion with the special places included (\code{brSpecialTorsionOrder})}
\addcontentsline{toc}{section}{E3. Torsion with special places}

\paragraph{Setting.} A component $cY$ lies on the curve $\mathcal C: Y^m=G(\theta)$
over $k(\theta)$, where $k=\Q(x)(t_1,\dots,t_r)$ and $\theta$ is the top monomial.
Its log part needs the order of the class of
\[ \delta=\sum_i b_i\,\Delta(\alpha_i)+\sum_l b'_l\,P_l \]
in $\operatorname{Pic}^0(\mathcal C)$. Here the $\Delta(\alpha_i)$ are the finite
residue divisors, and the $P_l$ are the places over the special points of
$\theta$: $\infty$ for a logarithm, $0$ and $\infty$ for an exponential
([B] Def.~2.21). Their residues $\tau_l$ enter the $\Q$-basis with the finite
ones. In genus~0 every degree-0 divisor is principal (v0.6533). In positive
genus, [B] gives no algorithm for the order of a divisor with special places;
this section supplies one.

\paragraph{The frame.} All divisors are written on the frame $\theta=a+1/u$
(\code{brShiftFrame}), with $a\in\Z$ ordinary. On this frame both special points
are finite: $\theta=\infty$ is $u=0$, and $\theta=0$ is $u=-1/a$. A special
place is then an ordinary finite place. It is unramified ($Y=\rho\,u_0$, with
$\rho^m=\kappa$) or totally ramified, and its prime ideal is built by
\code{brPointPlaceIdeal} or \code{placeIdealAt}. Every construction is generic
in the coefficient field, so the same code builds $\delta$ over $k$, over
$\Q$, or over $\mathrm{GF}(p)$.

\begin{definition}[good specialisation]
A point $\sigma\in\Q^{r+1}$ for $(x,t_1,\dots,t_r)$ is \emph{good} if it keeps
all of the following:
\begin{enumerate}
\item the degree and squarefree profile of $G$, hence the genus and the shape of
  the integral basis;
\item the degree of the branch-free denominator $B_p$, its squarefreeness, and
  its coprimality with $G$ and with $E'$;
\item the degree of every finite residue place $\Delta(\alpha_i)$
  (\code{brGoodSpecializationQ});
\item at each special location, a finite nonzero frame value $u_0$; at
  $\theta=0$, the multiplicity of $G$, and $B_p(0)\neq0$ when it holds
  generically (\code{brSpecialLocationsGoodQ});
\item the frame point $\theta=a$ ($u=\infty$) stays ordinary: $G_\sigma(a)$,
  $E_\sigma(a)\neq0$, and every polynomial of the frame ($G_u$, $B_{p,u}$, and
  the residue data $N_u$, $E_u$) keeps its $u$-degree
  (\code{brFrameDegreesKeptQ}). Without this a finite residue place can move to
  $u=\infty$ and leave the divisor.
\end{enumerate}
\end{definition}

\begin{theorem}\label{thm:E3}
Let $\sigma$ be good, and let $\delta_\sigma$ be the divisor built from the
specialised data on $\mathcal C_\sigma: Y^m=G_\sigma$ over $\Q$.
\begin{enumerate}
\item If $\delta_\sigma$ has infinite order, so has $\delta$.
\item If $\delta$ has finite order $N$, then $\delta_\sigma$ has order exactly $N$.
\end{enumerate}
\end{theorem}

\begin{proof}
Conditions (1)--(2) say that the plane model has the same singularity pattern
along the fibres over a neighbourhood $U$ of $\sigma$. The integral basis
specialises to the integral basis, so the normalisations form a smooth proper
family $\mathcal X\to U$ whose fibre over $\sigma$ is the normalisation of
$\mathcal C_\sigma$.

Conditions (3)--(4) say that each place in the support of $\delta$ specialises
to a place, or a set of places, of the same degree and multiplicity, and that
distinct places stay distinct. The special places stay over the same special
point and keep their ramification, and $u_0\neq0$ keeps the $m$ unramified
places apart. So $\delta$ extends to a relative divisor on $\mathcal X$ whose
fibre at $\sigma$ is $\delta_\sigma$. Equivalently, the ideal of $\delta$ built
over $k$ specialises to the ideal built from specialised data.

Specialisation $\operatorname{Pic}^0(\mathcal X_{\bar\eta})\to\operatorname{Pic}^0(\mathcal X_\sigma)$
is a group homomorphism, and in characteristic~0 it is injective on torsion:
for every $N$ the $N$-torsion of the relative Jacobian is a finite étale group
scheme over $U$, so its fibres all have the same order $N^{2g}$ and
specialisation is a bijection on them. Part~(2) follows, since the order
cannot drop. For part~(1), a relation $N\delta=\operatorname{div}(f)$ over
$k$ would specialise to $N\delta_\sigma=\operatorname{div}(f_\sigma)$.
\end{proof}

\paragraph{The order of $\delta_\sigma$.} We use Trager's two-prime method
([T] Ch.~6 \S2--3; \code{torsionOrderWith}) on $\mathcal C_\sigma$ over the
residue field $F_0$. A builder reconstructs $\delta_\sigma$ modulo $p$ with the
same field-generic code. A prime $p$ is good when Trager's conditions hold
(\code{goodPrime}) and also, for each special place (the \code{goodQ} hook):
\begin{itemize}
\item the frame value $u_0$ and the point $u=\mathrm{pt}$ are $p$-integral with
  $u_0\not\equiv0$;
\item $B_p(\mathrm{pt})\not\equiv0$, so the place does not collide with a finite
  residue place;
\item for an unramified place, $G_u(\mathrm{pt})\not\equiv0$ and the roots $\rho$
  stay distinct and nonzero;
\item for a ramified place, the multiplicity of $G_u$ at $\mathrm{pt}$ is kept.
\end{itemize}
Each of these is the reduction of a characteristic-0 non-vanishing, which is
checked first (\code{brSpecialPlaceDegenerate}). So all but finitely many primes
are good, and the prime search ends.

Under these conditions $\delta_\sigma\bmod p$ is the reduction of
$\delta_\sigma$. Trager's theorem (injectivity on prime-to-$p$ torsion) then
gives the only possible order $N$, or a certificate of infinite order when the
prime-to-$pq$ parts disagree.

\paragraph{Safeguards.}
\begin{itemize}
\item Before any order is computed, the specialised divisor is checked to have
  degree~0 on the frame: $\deg_u\det M-m\deg_u(\mathrm{den})=0$
  (\code{brFrameIdealDegree}). A lost or moved place is then an Internal error,
  never a certificate.
\item Trager's reduced arithmetic needs a root of $G_u$ of multiplicity prime to
  $m$. When characteristic~0 has none, no prime has one either, so the site stays
  \code{TorsionModP:frame-root} instead of an endless prime search.
\item Every characteristic-0 non-vanishing behind \code{goodQ} is checked first
  (\code{brSpecialPlaceDegenerate}).
\end{itemize}

\paragraph{Review (2026-11-14).} The independent review found a false
certificate in the first version. That version lacked condition~(5): for
$\frac{d}{dx}\log\bigl((L^3+L-x+59/5)^{1/3}-L\bigr)$ with $L=\log x$, the first
candidate $x\mapsto64/5$ put the residue place $\theta=x-59/5$ on the frame
point $\theta=a=1$. The fix is condition~(5) and the degree guard, and the
example is now a regression test. The review confirmed the rest:
\begin{itemize}
\item the theorem, given the corrected conditions;
\item leaving the $\tau_l$ out of Trager's distinct-residue condition (the
  divisor is defined by places);
\item the decision step.
\end{itemize}

\paragraph{Decision.} With $N$ in hand, $N\delta$ is tested for principality
\emph{generically}, over $k$ (\code{principalGenerator}).
\begin{itemize}
\item If $N\delta=\operatorname{div}(u)$, the log term is
  $(\beta_j/N)(\log u-\tfrac1m\log\operatorname{Norm}u)$.
\item If not, $\delta$ is not torsion. Indeed, a torsion $\delta$ has order
  exactly $N$ by Theorem~\ref{thm:E3}(2). The integral is then not elementary
  ([B] Prop.~2.13, [T] Ch.~6).
\end{itemize}

\paragraph{Tests.} \code{test\_special\_genus} in
\code{tests/test\_bronstein\_phase3.c} has direct calls of every routine, six
elementary integrals on genus-1 curves, and one certificate. Before v0.6548 all
of them stopped at \code{LogPart:special-genus}. In a 36-integrand sweep of
$m\ge3$ components, that site was the most frequent gap (6 of 36).

\section*{E6a. Reducible logarithms (\code{brReducibleLogSplit}, \code{brIntegrateKappaSplit})}
\addcontentsline{toc}{section}{E6a. Reducible logarithms}

\paragraph{Setting.} $\log a$ is reducible over the tower $k$: $D\log a=DL_0$
with $L_0=\sum r_i t_i+\sum s_j b_j$, where the $t_i=\log a_i$ are the
logarithmic monomials, the $b_j$ the arguments of the exponential monomials
$\theta_j=e^{b_j}$, and $r_i,s_j\in\Q$ (\code{Risch`LogReducible}). The rewrite
$\log a=L_0$ is wrong off a region: $\log x^2\ne2\log x$ for $x<0$. All
logarithms and powers are principal.

\begin{lemma}\label{lem:kappa}
The function $\kappa=\log a-L_0$ is locally constant. On each connected region
of analyticity its value $\kappa_0$ is either $0$ or transcendental over $\Q$.
\end{lemma}

\begin{proof}
$D\kappa=0$ off the branch cuts, so $\kappa$ is locally constant. Let $d\ge1$
be an integer with $dr_i,ds_j\in\Z$. Consider
\[ u=\frac{a^d}{\prod_i a_i^{dr_i}\prod_j\theta_j^{ds_j}}\in k. \]
The powers are integer powers, so $u$ lies in $k$. Moreover
$Du/u=d\,(Da/a-\sum r_iDa_i/a_i-\sum s_jDb_j)=d\,D\kappa=0$, so $u$ is a
constant of $k$, a nonzero element of $\Q$ (or of $\Qbar$ with algebraic
constants). As functions, $e^{d\log a}=a^d$, $e^{dr_i\log a_i}=a_i^{dr_i}$ and
$e^{ds_jb_j}=\theta_j^{ds_j}$, so $e^{d\kappa}=u$ pointwise. On a region,
$\lambda=d\kappa_0$ therefore satisfies $e^\lambda=u$ with $u$ algebraic. By the
Hermite--Lindemann theorem, $e^\lambda$ is transcendental for every algebraic
$\lambda\ne0$. So $\lambda$ is either $0$ or transcendental, and so is
$\kappa_0=\lambda/d$.
\end{proof}

Let $k'$ be the tower \emph{without} $\log a$: the monomials $t_i,\theta_j$ in
$L_0$ and everything else the integrand needs, but not $\log a$. We assume that
no monomial of $k'$ and no radicand depends on $\log a$, so $\log\log x^2$,
$e^{(\log x^2)^2}$ and $\sqrt{\log x^2}$ are all excluded. The field generated by
$k'$ and $\log a$ is then $k'(\kappa)$. Let $y$ be the radical, if any. Let
$f\in k'(y)[\log a]$ have degree $m$ in $\log a$. Expand it as
$f=\sum_{j=0}^m\kappa^jg_j$, with every $g_j\in k'(y)$. This is the substitution
$\log a\mapsto L_0+K$ followed by collecting powers of $K$.

\begin{theorem}\label{thm:kappa}
\leavevmode
\begin{enumerate}
\item If every $g_j$ has an elementary integral $G_j$, then
  $F=\sum\kappa^jG_j$ is an antiderivative of $f$ on every region.
\item On a region where $\kappa_0=0$, $f=g_0$, so $f$ is elementary there if
  and only if $g_0$ is.
\item On a region where $\kappa_0\ne0$, $f$ is elementary there if and only if
  every $g_j$ is.
\end{enumerate}
\end{theorem}

\begin{proof}
Part (1) holds because $D\kappa=0$, and (2) is immediate.

For (3), write $k$ for $k'(y)$ from here on. By Lemma~\ref{lem:kappa},
$\kappa_0$ is transcendental over $\Q$. A constant algebraic over $k$ is
algebraic over the constants of $k$ ([SI, Cor.~3.3.1]), so $\kappa_0$ is also
transcendental over $k$. (That is false over the tower \emph{with} $\log a$,
which contains $\kappa$; this is why $k'$ is needed.) Then
$k(\kappa_0)\cong k(K)$, with $K$ a new constant indeterminate, and the
isomorphism maps $f_K=\sum K^jg_j$ to $f$.

Suppose $f$ has an elementary integral on the region. By the strong form of
Liouville's theorem [R], [SI, Thm.~5.5.3],
\[ f_K=v'+\sum_{i=1}^n c_i\frac{u_i'}{u_i}, \]
with $c_i$ algebraic over $\Q(K)$, $v\in k(K)$, and
$u_i\in k(K)(c_1,\dots,c_n)^*$.

Now specialise. Let $\gamma$ be a primitive element of
$\Q(K)(c_1,\dots,c_n)$ over $\Q(K)$, with minimal polynomial $P(K,Z)$, taken
integral over $\Q[K]$. Then $P$ stays irreducible over $k(K)$, because in
characteristic~0 the constant field of $k$ is algebraically closed in $k$
([SI, Cor.~3.3.1]). So $L:=k(K)(c_1,\dots,c_n)=k(K)[Z]/(P)$, with basis
$1,\gamma,\dots,\gamma^{e-1}$ over $k(K)$.

For $q\in\Q$ outside the finitely many zeros of the discriminant of $P$ in $K$,
pick a simple root $\gamma_q\in\Qbar$ of $P(q,Z)$. Let $\sigma_q$ be the map
that sends $K\mapsto q$ and $\gamma\mapsto\gamma_q$. It is a ring homomorphism
from the subring of elements whose basis coefficients have no pole at $K=q$
into $k(\Qbar)$. Avoiding finitely many more $q$ puts $v$, the $u_i$, the
$1/u_i$ and the $c_i$ in that subring.

$\sigma_q$ commutes with $D$: $DK=0$, $D\gamma=0$, and $D$ acts on the basis
coefficients, elements of $k(K)$, coefficientwise in $K$. Applying it gives
\[ f_q=\sum q^jg_j=\sigma_q(v)'+\sum\sigma_q(c_i)\frac{\sigma_q(u_i)'}{\sigma_q(u_i)}. \]
So $f_q$ has an elementary integral over $k(\Qbar)$, and therefore over $k$:
the converse of Liouville's theorem needs only that the right-hand side is a
derivative of an elementary expression.

This holds for infinitely many $q$. Pick $m+1$ distinct ones,
$q_0,\dots,q_m$. The Vandermonde matrix $(q_l^j)$ is invertible, so each $g_j$
is a $\Q$-linear combination of $f_{q_0},\dots,f_{q_m}$, and hence elementary.

The converse of (3) is (1).
\end{proof}

\paragraph{The integrator's own expansion (v0.6558).}
The transcendental integrator underneath normalises every logarithm by
expanding its argument, $\log(ab)\mapsto\log a+\log b$ and
$\log(b^p)\mapsto p\log b$ (\code{rt\_expand\_logs}). The two sides have the
same derivative, but they differ by a locally constant $\kappa$ that does
\emph{not} drop out: it multiplies the other terms. For example
$\kappa=\log x^2-2\log x$ is $-2\pi i$ on $\Re x<0$, and in
$\int\bigl(\log x^2/(1+x)+2\log(1+x)/x\bigr)\,dx$ it contributes
$\kappa\log(1+x)$, not a constant. The integrator returned
$2\log x\log(1+x)$, which is wrong off $\Re x>0$.

The fix is the theorem above, applied to the expansion itself. Before the
tower is built, the $\kappa$-mode replaces every $\log a$ whose expansion
$L_0$ differs from it by $L_0+K$ and records $\kappa=\log a-L_0$. When $\kappa$
is identically $0$, $\log a$ is replaced by $L_0$ alone. That happens for a
positive rational factor ($\log(cu)=\log c+\log u$ off the cut) and for a
kernel that is already in the tower and is met again on the surface. So the
integrator only ever sees expanded logarithms, and Theorem~E6a decides the
coefficients as before. The answer above is $\log x^2\log(1+x)$; see
examples E6a-5 and E6a-6.

\paragraph{Verification on a cut reached through complex arithmetic.}
The verifier samples the real line as well, because an antiderivative must also
hold on a real-axis cut where the integrand is defined. A TrigToExp radicand
such as $i(e^{-ix}-e^{ix})/(e^{-ix}+e^{ix})=\tan x$ is real at real $x$, but it
is \emph{computed} through complex arithmetic. On the cut ($\tan x<0$), its
imaginary part is then rounding noise of either sign at every precision, so
the sample's value is undetermined. \code{PossibleZeroQ} now treats a
non-zero reading at such a point as uninformative and re-draws it. The test:
some radicand or log argument $b$ has $\Re b<0$ and $0<|\Im b|<2^{-150}|\Re b|$
at 200 bits. A radicand computed \emph{exactly} real, such as $(1-x)/(1+x)$,
is still sampled on its cut. Example E6d-4 is
$\int(\sqrt{\tan x}+\sqrt{\cot x})\,dx$.

\paragraph{What a certificate means.} A certificate here says that no
elementary expression is an antiderivative of $f$ on the whole domain. That is
weaker than it may look. $(\log x^2-2\log x)\sqrt{\log x}$ is certified
non-elementary, yet it vanishes identically on $\operatorname{Re}x>0$.
Theorem~\ref{thm:kappa} makes the statement exact region by region.

\paragraph{Decision, as implemented.}
\begin{itemize}
\item There is no ``$\kappa\equiv0$'' shortcut. A zero test cannot prove it:
  \code{PossibleZeroQ[Log[E\^{}x] - x]} is \code{True}, yet the difference is
  $-2\pi i$ wherever $\operatorname{Im}x\in(\pi,3\pi]$, as found in review.
\item Every $g_j$ is integrated. If all are elementary, the answer is
  $\sum\kappa^jG_j$, verified.
\item A non-elementary $g_j$ is a certificate when there is exactly one
  reducible logarithm and $\kappa$ is numerically nonzero at a sample point
  (\code{brNonzeroSomewhereQ}: $|\kappa(x_0)|>1/10$ at 30 digits, with sample
  points that include large imaginary parts). By
  Lemma~\ref{lem:kappa} the values of $\kappa$ form a discrete set, so a value
  that large is not a rounding artefact, and Theorem~\ref{thm:kappa}(3) applies
  on that region.

  If $g_0$ is the non-elementary one, (2) and (3) together certify $f$ on every
  region, without the sample.
\item A monomial or radicand that depends on $\log a$ violates the hypothesis on
  $k'$, and the split declines (\code{brReducibleLogSplit} returns
  \code{None}).
\item Two or more reducible logarithms stay \code{Tower:reducible-log}: their
  $\kappa_0$ may be algebraically dependent, so the specialisation argument does
  not apply as stated. So does $K$ in a denominator, which makes $K$ a symbolic
  parameter (Phase~4, E4).
\end{itemize}

\paragraph{Review (2026-11-14).} The independent adversarial review found no
false certificate in about 30 probes. It found an unstated hypothesis (the field
$k'$, now stated), a loose specialisation step (now done with a primitive
element), and an unsound shortcut (``$\kappa\equiv0$'' by a zero test, now
removed). It also found pre-existing lower-layer defects, recorded in
\code{tasks/todo.md}: the false zero above, an \code{Indeterminate} result of the
Automatic cascade for $(\log x^2-2\log x)e^{x^2}$, and an uncaught \code{Throw}
for $\int e^{x+\log(x^2)/4}$.

\paragraph{Tests.} \code{test\_reducible\_log} in
\code{tests/test\_bronstein\_phase3.c}. The certificate for
$(\log x^2-2\log x)\sqrt{\log x}$ uses Theorem~\ref{thm:kappa}(3) alone, with
$g_0=0$ and $g_1=\sqrt{\log x}$.

\section*{E6b. A cyclic radical compositum as one radical (\code{brCompositumPlan})}
\addcontentsline{toc}{section}{E6b. Cyclic radical compositum}

\paragraph{Setting.} We have radicals $\rho_i=b_i^{1/d_i}$ with
$b_i\in\Q(\mathbf v)$, whose Kummer group $\Gamma\subset(\Q/\Z)^J$ is cyclic
of order $N$ (\code{brKummerCyclicQ}). Here $\mathbf v$ is $x$ and the tower
symbols, $v_i\in(\Q/\Z)^J$ is the exponent vector of $\rho_i$ over the
irreducible factors of the $b_i$, and constant factors are ignored. As a field,
$\Q(\mathbf v)(\rho_1,\dots)$ is then one simple radical extension. Pointwise
presentations are another matter: a radicand $B$ with $B^{1/N}$ a generator need
not satisfy $\rho_j=R_jB^{k_j/N}$ pointwise.

\begin{proposition}
Suppose integers $c_i\in[0,d_i)$ and $k_j$ satisfy
\begin{enumerate}
\item[(a)] $c_iN/d_i\in\Z$ for all $i$;
\item[(b)] $(c_ik_j-\delta_{ij})/d_i\in\Z$ for all $i,j$;
\item[(c)] the class $\sum c_iv_i$ has order $N$ in $\Gamma$.
\end{enumerate}
Let $y=\prod\rho_i^{c_i}$. Then the identities
\[ y^N=B:=\prod b_i^{c_iN/d_i},\qquad \rho_j=R_j\,y^{k_j},\quad
   R_j=\prod_i b_i^{-(c_ik_j-\delta_{ij})/d_i}, \]
hold \emph{pointwise}, $B$ and the $R_j$ are rational, and $Y^N-B$ is
irreducible over $\Q(\mathbf v)$.
\end{proposition}

\begin{proof}
Only one identity is used: $(z^{1/d})^m=z^{m/d}$ for integer $m$, which holds
for the principal power because both sides are $\exp(\tfrac md\log z)$. So
$y^N=\prod_i\rho_i^{c_iN}=\prod_ib_i^{c_iN/d_i}$, an integer power of each
$b_i$ by (a). Likewise
$y^{k_j}/\rho_j=\prod_i\rho_i^{c_ik_j-\delta_{ij}}=\prod_ib_i^{(c_ik_j-\delta_{ij})/d_i}$,
an integer power of each $b_i$ by (b).

For irreducibility: by (c), the class of $B$ modulo $N$-th powers (constants
ignored) has order $N$. So $B$ is not a $p$-th power for any prime $p\mid N$,
and when $4\mid N$ it is not of the form $-4w^4$. By Capelli's theorem,
$Y^N-B$ is irreducible.
\end{proof}

\paragraph{Use.}
\begin{itemize}
\item The integrand is integrated over $\Q(\mathbf v)(Y)$, $Y^N=B$, with $Y$
  written as $B^{1/N}$.
\item The answer is mapped back by $B^r\mapsto y^{rN}$ (\code{brCompositumBack}).
  $Y\mapsto y$ is an embedding of differential fields: both satisfy $Y^N=B$ and
  therefore $DY=Y\,DB/(NB)$. So the image of an antiderivative is an
  antiderivative, and a certificate transfers.
\item When no $c,k$ exist, the site stays \code{Tower:cyclic-compositum}. One
  example is $\sqrt{x^2(x+1)}$ with $\sqrt{x+1}$, which agree only up to a
  locally constant sign. Another is $\sqrt{2x}$ with $\sqrt x$, which needs a
  constant field extension.
\end{itemize}
The search is over $\prod d_i\le20000$ tuples. Condition (c) follows from (b):
$v_j=k_j\sum_ic_iv_i$, so $\sum c_iv_i$ generates $\Gamma$. It is kept as a
cheap filter.

The back-map is safe. It acts only on powers $B^r$ whose $rN$ is an integer, and
anything it leaves (a CAS that split a constant out of $B^{1/N}$) is caught by
the final verification against the user's integrand.

\paragraph{Tests.} \code{test\_compositum\_plan},
\code{test\_compositum\_rewrite} and \code{test\_compositum\_integrate} in
\code{tests/test\_bronstein\_phase3.c}. They include a numeric pointwise check of
the identities off the real axis, and the Chebyshev certificate for
$\sqrt x\,(1+x)^{1/3}$.

\section*{E6d. One radical up to a locally constant sign (\code{brIntegrateSignSplit})}
\addcontentsline{toc}{section}{E6d. Up to a sign}

\paragraph{Setting.} This is E6b when the conditions there fail. The usual reason
is a radicand with a square factor: $\sqrt{x^2(x+1)}$ and $\sqrt{x+1}$ generate
the cyclic group $\Z/2$, yet $\sqrt{x^2(x+1)}=\pm x\sqrt{x+1}$ only up to a sign
that changes between regions.

\begin{proposition}
Suppose $c,k$ satisfy (a) and (c) of E6b, and (b) is weakened to
\begin{enumerate}
\item[(b$'$)] $e_{ij}=(c_ik_j-\delta_{ij})/d_i\in\tfrac12\Z$, and for each $j$
  the product $\prod_{i\in H_j}b_i$ over $H_j=\{i:e_{ij}\notin\Z\}$ equals
  $s_j^2$ with $s_j\in\Q(\mathbf v)$.
\end{enumerate}
Then pointwise
\[ \rho_j=\varepsilon_j\,R_j\,y^{k_j},\qquad
   R_j=\Bigl(s_j\prod_ib_i^{\lfloor e_{ij}\rfloor}\Bigr)^{-1},\qquad
   \varepsilon_j=\prod_{i\in H_j}b_i^{1/2}\big/s_j, \]
and $\varepsilon_j$ is locally constant with values in $\{1,-1\}$.
\end{proposition}

\begin{proof}
As in E6b, $y^{k_j}/\rho_j=\prod_ib_i^{e_{ij}}$ pointwise. Also
$b^{n+1/2}=b^n b^{1/2}$ for the principal power. Hence
$y^{k_j}/\rho_j=\prod_ib_i^{\lfloor e_{ij}\rfloor}\prod_{H_j}b_i^{1/2}$. The
last product squares to $s_j^2$, so it equals $\varepsilon_js_j$ with
$\varepsilon_j^2=1$. Being continuous wherever the radicals are analytic,
$\varepsilon_j$ is locally constant.
\end{proof}

A constant factor must be a positive rational square
(\code{brRationalSquareRoot}). Otherwise $\varepsilon_j$ would take values in
$\{\pm\sqrt c\}$ or $\{\pm i\}$, which is a constant field extension (Phase~4),
and so would roots of unity of order greater than~2.

\paragraph{Decision.} Write $\varepsilon_j$ as a symbol $E_j$, and for each sign
pattern $s\in\{\pm1\}^J$ let $f_s=g|_{E=s}$, an element of one simple radical
extension. Distinct pairs $(H_j,s_j)$ give distinct symbols; the patterns that
never occur are harmless for answers.
\begin{itemize}
\item If every $f_s$ has an elementary integral $F_s$, then
  $F=\sum_s\prod_j\tfrac12(1+s_j\varepsilon_j)\,F_s$ is an antiderivative on
  every region. The product is the indicator of the region where $\varepsilon=s$,
  and its derivative is~0.
\item If some $f_s$ is not elementary and the pattern $s$ is taken on an
  \emph{open} region, then $f=f_s$ there, and $f$ has no elementary
  antiderivative on the whole domain. The certificate means what it means in
  E6a.

  The sampler (\code{brSignPatternSeenQ}) accepts a point only if every
  radicand of $\varepsilon$ is off its cut $(-\infty,0]$ there, and the same
  values are taken at four perturbations of size $10^{-6}$. On a cut, a product
  of principal radicals can take a value that no open set has:
  $\sqrt{x+1}\sqrt{1/(x+1)}=-1$ exactly for real $x<-1$.
\item An unseen pattern with a non-elementary $f_s$ stays
  \code{Tower:cyclic-compositum}.
\end{itemize}

\paragraph{Review (2026-11-14).} The first release (v0.6549) sampled real points
on those cuts. That gave a false certificate for
$\tfrac12(1+e^{x^2})+\tfrac12(1-e^{x^2})\sqrt{x+1}\sqrt{1/(x+1)}$, which equals~1.
v0.6550 fixes the sampler, and the example is a regression test.

\paragraph{Tests.} \code{test\_sign\_compositum} in
\code{tests/test\_bronstein\_phase3.c}, and examples E6d-1 to E6d-4 in the gap log.

\section*{E6e. The log part has no coordinates of weight $k\ne1$ (gap \#12)}
\addcontentsline{toc}{section}{E6e. Weights of the log part}

\paragraph{Setting.} A component $cY$ lies on $\mathcal C: Y^m=G$ over
$k(\theta)$, $m\ge3$. Let $\sigma: Y\mapsto\zeta Y$, with $\zeta$ a primitive
$m$-th root of unity, after the constants are extended by $\zeta$. On the basis
$w_i=Y^i/d_i$, an element $v$ has coordinates $v_i$, and $\sigma$ acts on them
by $\zeta^i$. The log part is $L=\sum_j(\beta_j/N_j)\log u_j$ with
$N_j\delta_j=\operatorname{div}(u_j)$. The implementation keeps its weight-1
projection through the norm formula
$L_1=\sum(\beta_j/N_j)(\log u_j-\tfrac1m\log\operatorname{Norm}u_j)$. That
formula is valid exactly when the coordinates $2,\dots,m-1$ of $L'$ vanish,
which [B] and [T] state only for $m=2$. Since v0.6533 a nonzero coordinate was
reported as \code{LogPart:weights}. It was never observed. A sweep of 120
random components with $m\in\{3,4\}$ (logarithms and exponentials) and the 36
integrands of E3 did not trigger it. The following explains why.

\begin{theorem}\label{thm:E6e}
Suppose the residue function $P\mapsto r(P)=\operatorname{res}_P(L')$ on the
places of $\mathcal C$ is a sum $r=r_1+r_0$ of a $\sigma$-equivariant part of
weight one, $r_1(\sigma^{-1}P)=\zeta\,r_1(P)$, and a $\sigma$-invariant part,
$r_0(\sigma P)=r_0(P)$. Suppose also that the curve is reduced (the constants
of $K$ algebraic over $k$ lie in $k$) and that $\sigma$ commutes with $D$. Then
for every $k\not\equiv0,1\pmod m$ the weight-$k$ part
$L_k=\tfrac1m\sum_{j=0}^{m-1}\zeta^{-kj}\sigma^j(L)$ has $L_k'=0$, so the
coordinates $2,\dots,m-1$ of $L'$ vanish. ($L_0'$ lies in $k$ and need not be
$0$: scaling a $u_j$ by $v\in k$ adds a multiple of $v'/v$.)
\end{theorem}

\begin{proof}
Fix the convention $\operatorname{res}_P(\sigma^j g)=\operatorname{res}_{\sigma^{-j}P}(g)$
($\sigma^j g$ is $g\circ\sigma^{-j}$ on places), which holds because $\sigma$
commutes with $D$. Iterating the hypothesis gives
$r_1(\sigma^{-j}P)=\zeta^jr_1(P)$. Hence
\[ \operatorname{res}_P(L_k')=\tfrac1m\sum_j\zeta^{-kj}\bigl(\zeta^{j}r_1(P)+r_0(P)\bigr)
   =r_1(P)\tfrac1m\sum_j\zeta^{(1-k)j}+r_0(P)\tfrac1m\sum_j\zeta^{-kj}, \]
which is $0$ for $k\not\equiv0,1$. The residue of $cY$ is equivariant in exactly
this way, since $\sigma(cY)=\zeta cY$. A place fixed by $\sigma$, or with a
nontrivial stabiliser (partial ramification when $m=4$), gets residue~0 from
the same identity. So $L_k'$ is a combination $\sum_ic_i\,u_i'/u_i$, with
$c_i\in\Q(\zeta,\beta)$, that has zero residue at every place.

Choose a $\Q$-basis $e_1,\dots,e_r$ of the span of the $c_i$, and write
$\sum_ic_iu_i'/u_i=\sum_le_l\,v_l'/v_l$ with $v_l$ products of integer powers
of the $u_i$. The residue at $P$ is $\sum_le_l\operatorname{ord}_P(v_l)=0$, with
integer orders, so every $\operatorname{ord}_Pv_l=0$. Each $v_l$ therefore has
divisor zero and lies in the algebraic closure of $k$ in $K$, which is $k$
because the curve is reduced. So $v_l'/v_l\in k$ is fixed by $\sigma$, and has
weight~0.

For $k\not\equiv0$ a weight-0 element has no weight-$k$ part, so $L_k'=0$.
\end{proof}

\paragraph{Where the hypotheses hold.}
\begin{itemize}
\item \textbf{Finite places, special places of a logarithm.} The residue function
  of the log terms is $r_1+r_0$. Here $r_1$ is the residue of $cY$, and $r_0$
  is $\sigma$-invariant: it comes from the reference divisor
  $\kappa(\theta-c)_0$ of the orbit construction, and from the places over
  $\theta=0$ of an exponential, where $cY$ may have ordinary residues that the
  log part does not. All of this holds by construction: $\sum_j\beta_j\delta_j=\sum_i\alpha_i\Delta(\alpha_i)$
  (and the special taus). Generator normalisation multiplies $u_j$ by an
  element of $k(\theta)$, which is $\sigma$-invariant and changes only
  weight~0. A logarithm's special residues are ordinary residues of
  $cY\,d\theta/\eta$.
\item \textbf{Special places of an exponential.} Bronstein's residues there
  ([B] Def.~2.21, with the reference value $c_{\bar p}$) are not ordinary
  residues of a differential. The review suggests that
  $\tau_p=\varepsilon r(c_p-c_{\bar p})$ also splits into a weight-1 part and an
  invariant part. Until that is written out, this path keeps
  \code{LogPart:weights}.
\end{itemize}

\paragraph{Implementation.} On every other path a nonzero coordinate is an
Internal invariant violation, \code{LogPartWeightNonzero}: it can be a bug, but
not a missing branch.

\paragraph{Review (2026-11-14).} The independent review confirmed the result.
It corrected three things in the first write-up, all fixed above:
\begin{itemize}
\item an inconsistent residue convention;
\item the claim for weight~0;
\item a hypothesis stronger than what holds at the places over $\theta=0$.
\end{itemize}
About 100 probes with $m\in\{3,4\}$ never raised \code{LogPartWeightNonzero}.

\paragraph{Tests.} \code{test\_weights} in \code{tests/test\_bronstein\_phase3.c}.

\section*{E6c. Exponents with a rational log-linear part (\code{brSplitLogExponent})}
\addcontentsline{toc}{section}{E6c. Log-linear exponents}

For rational $r_i$,
\[ e^{u+\sum r_i\log a_i}=e^u\prod a_i^{r_i} \]
holds pointwise: $\exp(v+w)=\exp(v)\exp(w)$, and $\exp(r\log a)=a^r$ is the
definition of the principal power. So $E^{x+\log(x)/2}$ is $\sqrt x\,e^x$.
Without the split it would be a new hyperexponential monomial over
$\Q(x)(\log x)$, whose Risch DE needs a radical the tower does not have.

The same identity with $u=0$ turns $E^{r\log a}$, found by
\code{Risch`ExpReducible} with a fractional coefficient on a logarithm, into
the radical $a^r$ (\code{brExpRewrite}).

A monomial left unused by the rewrite is pruned (\code{brPruneMonomials}), so
it cannot pose as a radical in the middle of the tower. For example,
$\log\log x$ is unused once $E^{x+\log(\log x)/2}=\sqrt{\log x}\,e^x$.

\paragraph{Tests.} \code{test\_exp\_of\_log} in
\code{tests/test\_bronstein\_phase3.c}.

\section*{E4.1. The constant field $\Q(i)$ (\code{tragerConstantI})}
\addcontentsline{toc}{section}{E4.1. Constant field Q(i)}

\paragraph{Setting.} Integrands whose constants are Gaussian rationals, such as
$\sqrt{x+i}/x$, $e^{ix}\sqrt{1+e^{ix}}$ and $\sqrt{\log x+i}$. The constant
field is $C=\Q(i)$, and $i$ stays the literal \code{I} of the CAS, whose
arithmetic reduces $i^2=-1$ in every \code{Together}, \code{Expand} and zero
test.

The algorithms of [B], [T] and [SI] hold over any constant field of
characteristic~0, with three provisos:
\begin{enumerate}
\item factorisations, place degrees and residue fields are taken over the
  constant field actually in use;
\item every residue field is a field (no zero divisors in its zero tests);
\item reduction modulo $p$ maps the constants consistently.
\end{enumerate}
The implementation meets each of these as follows.

\begin{proposition}\label{prop:E41}
With the flag \code{tragerConstantI} set, which it is exactly when the
integrand has a Gaussian rational:
\begin{enumerate}
\item every factorisation of a characteristic-0 polynomial is over $\Q(i)$
  (\code{tragerFactorList} adds \code{Extension -> I}). A relative extension
  $\Q(i)[t]/(\mu)$ built from an irreducible factor $\mu$ is therefore a field;
\item an absolute residue field $\Q[t]/(\mu)$ from \code{ToNumberField} is built
  from the residues \emph{and} $i$. It contains~$i$, and every Gaussian
  rational is written $a+b\,\mathrm{rep}(t)$ before reduction modulo~$\mu$
  (\code{fieldSubI} in \code{kc} and \code{rc}). So no element mixes a literal
  $i$ with $t$;
\item reduction modulo $p$ (torsion) refuses data with a literal $i$
  (\code{TorsionModP:complex}). Positive-genus torsion over $\Q(i)$ is therefore
  undecided, never wrongly decided;
\item every coefficient test accepts exactly the Gaussian rationals
  (\code{constantFieldNumberQ}, \code{brRationalInQ}).
\end{enumerate}
\end{proposition}

\begin{proof}[Why (2) is needed]
If $\mu\in\Q[t]$ is irreducible over $\Q$ but splits over $\Q(i)$ (as
$t^2+1$ does), then $\Q(i)[t]/(\mu)$ has zero divisors. An element
$a(t)+i\,b(t)$ could then be zero without $a=b=0$, so a zero test that treats
$i$ and $t$ independently is wrong. Adjoining $i$ to the absolute field and
substituting removes the independent $i$. Then $\Q[t]/(\mu)$ is a field
($\mu$ irreducible over $\Q$) containing $\Q(i)$, and its zero tests are exact.
\end{proof}

\paragraph{Spike and result.} The front end with only the declines lifted
answered 2 of 12 integrals in a $\Q(i)$ battery (\code{Q}-only tests rejected
$i$ at about eight points). With the layer above, 10 are answered and verified,
1 is certified, and 1 remains (a torsion question, refused mod $p$). The
certificate $\int\sqrt{\log x+i}\,dx$ was checked by hand: FriCAS reports
``implementation incomplete (constant residues)'' on it. With $u=\log x$ it is
$\int e^u\sqrt{u+i}\,du$, and the Risch DE $b'+b(1+\tfrac1{2(u+i)})=1$ has no
rational solution.

\paragraph{Review (2026-11-14): false certificates in v0.6551.} The independent
review found that v0.6551 certified elementary integrals such as
$\int\sqrt{e^x+i}\,dx$ non-elementary, among them $(e^x+i)^{1/3}$,
$\sqrt{1+i\log^2x}/x$ and $e^x/\sqrt{1+ie^{2x}}$.
\begin{itemize}
\item \textbf{Cause.} \code{brGammaForm} decided whether $\kappa$ in
  $\gamma=\kappa^{1/m}u$ is constant by \code{FreeQ[kappa, \_Symbol]}, which is
  false for $\kappa=i$ (the head \code{Complex} is a symbol). The constant $i$
  was taken for a radical $\gamma$, the special residue was set to~$0$ by a
  Galois-weight argument that needs $\gamma$ \emph{not} constant, and the Risch
  DE then ``proved'' that no algebraic part exists.
\item \textbf{Fix (v0.6552).} A constant is \code{NumericQ}, and a non-real
  Gaussian $\kappa$ is refused (\code{Residues:gaussian-kappa}) until special
  residues over $\Q(i)(\kappa^{1/m})$ are implemented.
\end{itemize}
The review found three further faults, all fixed in v0.6552:
\begin{itemize}
\item a RootSum term passed through the real-form presentation leaked a private
  symbol, and \code{PossibleZeroQ} accepted the result. Now \code{presentable}
  skips integrands with complex constants and \code{realForm} skips RootSums;
\item the tower branch of \code{residueField} ignored the flag. Now it refuses;
\item the mod-$p$ refusal checked a fixed list of keys. Now \code{gfMod}, the
  single choke point of every reduction, refuses complex data.
\end{itemize}
Proposition~\ref{prop:E41}(3) now holds at that choke point.

\paragraph{Next.} Map $i$ to a square root of $-1$ modulo $p\equiv1\pmod4$ in
the torsion code; then the trigonometric front end (\code{TrigToExp}); then
general number fields through \code{fieldFn[vars, \{alpha, mu, theta\}]}.

\paragraph{Tests.} \code{tests/test\_bronstein\_phase4.c}.

\section*{E4.1b. $i$ modulo $p$, Gaussian residues and $\kappa$}
\addcontentsline{toc}{section}{E4.1b. i modulo p}

\paragraph{Reduction.} Let $p\equiv1\pmod4$ and let $r$ be the least root of
$z^2+1$ modulo $p$ (\code{gfSqrtMinusOne}). Then $\mathfrak p=(p,i-r)$ is a
prime of $\Z[i]$ of degree one, with residue field $\mathrm{GF}(p)$, and
$a+bi\mapsto a+br$ is the reduction map. Trager's torsion theory ([T]
Ch.~6~\S2--3) holds for curves over a number field at a prime of good reduction
of degree one, and so does its specialisation step in E3.

Two conditions make the implementation that map:
\begin{enumerate}
\item every reduction goes through \code{gfMod}, which applies it;
\item when a residue field contains $i$ as $\mathrm{rep}(t)$, the root $\rho$ of
  $\mu$ modulo $p$ is chosen with $\mathrm{rep}(\rho)\equiv r$
  (\code{goodPrime}). Otherwise a literal $i$ and $\mathrm{rep}(t)$ would reduce
  through the two conjugate primes.
\end{enumerate}
Primes $p\not\equiv1\pmod4$ are skipped when the data carry~$i$; the good
primes keep a positive density. A literal $i$ modulo such a prime is refused,
never reduced.

\paragraph{Gaussian residue polynomials.} An irreducible factor $q$ over $\Q(i)$
with non-rational coefficients is handed to \code{ToNumberField} through the
roots of its norm $q\bar q\in\Q[z]$, picked numerically
(\code{gaussianRootObjects}). The field built contains those roots and $i$.

\paragraph{Gaussian $\kappa$.} Under $\Q(i)$, a Gaussian $\kappa$ at the special
places is a constant, and the special-place construction is unchanged. Its
residue field $\Q(i)(\text{roots of }z^m-\kappa)$ is built as above. The
false certificates of v0.6551 came from reading $\kappa=i$ as non-constant
(E4.1), not from this path.

\paragraph{Exact identity over a field containing $i$.} The component identity
check reduces modulo $\mu$ only after writing the literal $i$ in the generator
(\code{fieldSubI}). Before, a true identity read as nonzero, an Internal error
and not a wrong answer.

\section*{E4.2. The trigonometric front end (\code{brTrigFrontEnd})}
\addcontentsline{toc}{section}{E4.2. Trigonometric front end}

\code{TrigToExp} rewrites the circular and hyperbolic functions and their
inverses into exponentials and logarithms over $\Q(i)$. For example,
$\arcsin x=-i\log(ix+\sqrt{1-x^2})$ is a logarithm together with the radical.
That rewrite is an identity of the functions the CAS evaluates, so:
\begin{itemize}
\item the integrand $g$ is integrated over $\Q(i)$ (E4.1);
\item an antiderivative verified against $g$ is one of $f$;
\item a certificate for $g$ is one for $f$.
\end{itemize}

A second numeric check against $f$ was dropped. For
$\int\sqrt{\tan x}\,dx$ the answer carries Root objects of a quartic, and
\code{PossibleZeroQ} at complex points was numerically unstable on it: at one
sample it gave $D F-g=3i$ while $f-g=0$.

The trigonometric presentation \code{ExpToTrig[F]} is returned only when it is
itself verified against~$f$. A trigonometric integrand needs at most one radical
in the exponential form, and the scope test runs on that form
(\code{BronsteinScope}).

\paragraph{Results.} On a 13-integral battery, 9 were answered and verified,
including $\sqrt{\tan x}$, $\tan^{1/3}x$, $\sqrt{1+\sin x}$,
$\arcsin x/\sqrt{1-x^2}$ and $x\arctan x/\sqrt{1+x^2}$. The other 4 were
certified:
\begin{itemize}
\item $\sqrt{\sin x}$ and $\sqrt{\cos x}$, which are elliptic;
\item $\sec^{1/3}x$;
\item $\sqrt{\log x}\sin x$.
\end{itemize}
The answers for $\sqrt{\tan x}$ still carry Root objects of the residue field.
Rewriting them in radicals and in a real form is presentation work for later.

\paragraph{Review (2026-11-14).} The independent review ran about 150 probes and
found no false certificate and no wrong antiderivative. It confirmed that one
reduction map is used on every path (E4.1b), the soundness of a Gaussian
$\kappa$, and \code{TrigToExp} at 14 points per head.

Three of its findings were fixed in v0.6554:
\begin{itemize}
\item The hyperbolic presentation \code{ExpToTrig} wrote $e^{-kx}$ as
  $\cosh kx-\sinh kx$, which cancels catastrophically at machine precision. It is
  now only used for circular integrands, and never when it introduces a
  hyperbolic head.
\item Residues all in $\Q(i)$ took the $\Q$ shortcut of \code{residueField},
  which treated $\Q(i)$ as one-dimensional. The field is now built with $i$.
\item A Gaussian resultant was left by \code{Expand} as a sum of uncancelled
  fractions. \code{Together} is now applied first.
\end{itemize}

Two remarks:
\begin{itemize}
\item \textbf{The integrand's own branch cut.} At real points where the radicand
  of $f$ is negative, $f$ sits on its own branch cut. There the presented $F$
  can reach the other side of the cut, so its machine values differ from those
  of $f$, although $F'=f$ on every open set. This is why the antiderivative is
  verified with complex sampling.
\item \textbf{A latent issue with \code{ArcSech}.} Machine \code{ArcSech} at
  negative reals disagrees with its exact form and with \code{TrigToExp}.
  Integrands with \code{ArcSech} are declined today: the rewrite has two
  radicals. If that changes, the identity claim must be rechecked.
\end{itemize}

\paragraph{Tests.} \code{tests/test\_bronstein\_trig.c}, and gap-log entries
E41b-1 to E41b-3 and E42-1 to E42-5.

\section*{E4.3. Number-field constants (\code{tragerAlgConstants})}
\addcontentsline{toc}{section}{E4.3. Number-field constants}

\paragraph{Scope.} The constants of the integrand are radicals of rationals,
such as $\sqrt2$, $2^{1/3}$, $\sqrt5$, possibly together with~$i$. The constant
field is $C=\Q(i?,\text{constants})$, and the constants stay literal, as $i$
did in E4.1.

\paragraph{What carries over from E4.1.}
\begin{itemize}
\item Factorisations are over $C$ (\code{tragerFactorList} with
  \code{Extension}).
\item An absolute residue field is built with every constant adjoined, and
  records each constant's representation in its generator. The canonicalisers
  rewrite \emph{every} radical of a rational in the generator before reducing
  modulo~$\mu$ (\code{fieldSubI}, \code{fieldConstRep}). That includes the
  products the arithmetic forms, such as $\sqrt2\sqrt3\to\sqrt6$.
\item Residue polynomials with algebraic coefficients reach
  \code{ToNumberField} through their norm. The norm is a resultant against each
  constant's minimal polynomial, and the roots are picked numerically.
\item A radical of a constant inside a radicand ($\sqrt{x+\sqrt2}$) is a
  constant, not nesting.
\end{itemize}

\paragraph{What does not carry over: certificates.} Trager's component
integrator works over \code{fieldQ[]}, whose canonical form is
\code{Expand}. That is exact for Gaussian rationals, because the CAS
rationalises their denominators. It is not exact for radicals:
$1/(1+\sqrt2)$ is left as it is, so a zero coefficient can read as nonzero, and
a ``no'' is not a proof.

On a 14-integral battery this gave two false verdicts before the gate:
\begin{itemize}
\item the derivative of $\sqrt{x^3+\sqrt2}+\log(\sqrt{x^3+\sqrt2}+x)$;
\item $(1+\sqrt3)\sqrt{\tan x}$.
\end{itemize}
So with radical constants, a non-elementarity verdict is reported as
``implementation incomplete'' (\code{NF:certificate}). Answers are unaffected:
they are verified.

Linearity removes the constants where it can. A constant factor leaves the
integrand, $\int c\,g=c\int g$, so $\sqrt2\sqrt{\log x}$ is decided, with a
certificate, over~$\Q$. Phase~4.3b makes the constant field a proper field record,
\code{fieldNF[t, mu, gamma]}, in the base-field arithmetic.

\paragraph{Torsion.} Reduction modulo $p$ refuses radical constants
(\code{TorsionModP:algebraic-constant}). Phase~4.3b reduces at a degree-one
prime of $C$.

\paragraph{A bug caught on the way.} The first version returned the residue
field from inside a nested \code{Module}. In this CAS a \code{Return} there only
leaves the inner scope. The generic field, without its $i$ key, was returned
instead. That gave a false certificate on $\int\sqrt{\tan x}\,dx$, a
\emph{regression} of E4.2, caught by the battery before release.

\paragraph{Results.}
\begin{itemize}
\item 8 of 14 battery integrals are answered and verified, including
  $\sqrt{x+\sqrt2}/x$, $\sqrt{x+2^{1/3}}/x$ and $x/\sqrt{x^4+\sqrt5}$.
\item 2 are certified over $\Q$ by linearity.
\item The rest are honest declines, recorded in the gap register.
\end{itemize}

\paragraph{Review (2026-11-14).} The independent review ran about 940 probes and
found two false certificates, both fixed before release.
\begin{itemize}
\item \textbf{Linearity needs a provably nonzero factor.} The rule
  $\int cg=c\int g$ certified $c\,g$ with $c=(1+\sqrt2)^2-3-2\sqrt2=0$. It now
  applies only to an algebraic $c$, which \code{RootReduce} decides exactly; a
  zero $c$ gives~$0$. A transcendental $c$ (for example $\log4-2\log2$) cannot be
  decided this way, so it takes the normal path.
\item \textbf{ParallelMixedTower has the same arithmetic gap.} With radical
  constants its certificates are demoted too.
\end{itemize}
It also confirmed that every other verdict with radical constants was demoted,
and that a 110-probe sweep on $\Q$ and $\Q(i)$ held.

\paragraph{Tests.} \code{tests/test\_bronstein\_nf.c}, and gap-log entries
E43-1 to E43-6.

\section*{E4.3b. The constant field in the base arithmetic and at the primes (v0.6559)}
\addcontentsline{toc}{section}{E4.3b. Constant field in base arithmetic and mod p}

\paragraph{Base-field arithmetic.} With radical constants $c_1,\dots,c_r$ in the
integrand, the constant field is $K_0=\Q(c_1,\dots,c_r)$. Trager's component
pipeline (Hermite reduction, the branch-free denominator, the residue
polynomial) was written over $\Q$ and treated each $c_i$ as an independent
symbol. Two consequences turned elementary integrands into false verdicts,
both hidden by the demotion of E4.3:
\begin{enumerate}
\item The squarefree decomposition of $z^2(1-\sqrt[3]2\,z)$ was not split at
  all, so Hermite reduction stopped at once and the remainder of
  $\int(x-\sqrt[3]2)^{-2/3}\,dx$ was declared a holomorphic differential.
\item A combination of fractions whose common denominator needs
  $\sqrt2^{\,2}=2$ stayed an uncombined sum. Hermite stopped early, and a
  pole at infinity was reported for an elementary integrand on
  $y^2=x^3+\sqrt2$.
\end{enumerate}
Every rational-function operation on the base field now goes through
$K_0$-aware helpers: a squarefree list, \code{Together} and \code{PolynomialGCD}
with \code{Extension} $\to$ \code{Automatic} (\code{tragerSquareFreeList},
\code{tragerTogether}, \code{tragerGCD}), and the canonical form
\code{rc[fieldQ[], e]}. Over $\Q$, or $\Q(i)$ alone, the helpers are the plain
operations, so existing behaviour is unchanged.

\paragraph{Reduction at a prime of degree 1.} The torsion test of [T] reduces
the curve, the divisor and the residues modulo a prime. Over $K_0$ the
reduction must be one ring homomorphism $\mathcal O_{K_0}\to\mathbb{F}_p$ for all the
data. Let $\gamma$ be a primitive element of $K_0$ with monic minimal
polynomial $\mu$, and write each constant as $c_i=P_i(\gamma)$, $P_i\in\Q[g]$.
Let $p$ divide no denominator of $\mu$ or of any $P_i$ and not divide
$\operatorname{disc}\mu$, and let $r$ be a simple root of $\mu$ mod $p$. Then
$\gamma\mapsto r$ is the reduction at the prime $\mathfrak p=(p,\gamma-r)$ of
degree 1, and $c_i\mapsto P_i(r)$. Taking the \emph{least} such $r$ for every
$p$ makes all reductions at $p$ agree (\code{gfConstRoot},
\code{gfConstImage}), as \code{gfSqrtMinusOne} does for $\Q(i)$. When $I$ is
also in use it is included in $K_0$ and reduced the same way. A residue field
$F\supseteq K_0$ is reduced at a root $\rho$ of its own minimal polynomial.
\code{goodPrime} keeps only the $\rho$ under which the recorded representations
of the constants in $F$ reduce to the same images, so the reduction of $F$ lies
over $\mathfrak p$.

The reduction is injective on the torsion of order prime to $p$ at a prime of
good reduction of a curve over a number field [Silverman, VII.3.1]. This is
the $\mathfrak p$-adic form of the argument already used over $\Q$, so two good
primes with disagreeing orders still prove infinite order. With this,
$\int dx/\bigl((x-\sqrt2)\sqrt{x^3+1}\bigr)$ is decided: its orders modulo 17
and 71 disagree. Its certificate stays demoted until the residue and log
fields above $K_0$ are exact as well (steps C3 and C6) and the whole path has
had an adversarial review (C8).


\paragraph{Residue and log fields above $K_0$ (v0.6562).} A residue factor $q$
whose coefficients carry the constants has roots generating
$K_0(\alpha)$. The relative presentation $\Q[Z]/(q)$ treats the literal
constants as numbers, and $1/(1+\sqrt2)$ is never rationalised, so a zero can
read as non-zero. The field of one root is therefore built \emph{absolute},
$\Q(\alpha,c_1,\dots,c_r)=\Q(\delta)$ with $\delta$ a primitive element, and
the constants' representations are recorded so that every literal constant is
written in $\delta$ (\code{fieldOfRoot}). The same holds for the canonical form
over the base: while constants are in use, \code{kc[fieldQ[], e]} is the normal
form in the primitive element of $K_0$, with denominators rationalised by the
norm (\code{kcConstQ}).

The trace shortcut of [T] divides $\operatorname{Tr}_{\Q(\gamma)/\Q}$ by the
relative degree. That is valid when $\Q(\gamma)$ is an extension of $\Q$ of
that degree, not of $K_0$, so it gives no verdict over $K_0\ne\Q$, as already
over $\Q(i)$; the complete construction decides instead.


\paragraph{Bronstein's $k$-arithmetic over $K_0$ (v0.6564).} Hermite reduction
in $\theta$ needs the squarefree decomposition, exact quotients and the
extended gcd in $k[\theta]$ with $k=\Q(x,\text{lower})(K_0)$. The C kernel of
the decomposition loses a polynomial whose coefficients carry radicals (it
returns nothing), and this produced two false certificates (examples E43b-8 and
E43b-9). The decomposition is now [B, SplitSquarefreeFactor] over $K_0$, and
the quotient, the congruence and the cancellation of common $\theta$-factors
are computed in $\code{fieldFn}[\{x,\dots\},\{g,\mu,\gamma\}]$, exactly modulo
$\mu$.

\paragraph{Radicals equal over a larger constant field.} For $c>0$ rational and
any $u$, $(cu)^r=c^ru^r$ on every branch, because $\arg(cu)=\arg u$. Pulling the
positive content out of each radicand therefore turns $\sqrt{2x}+\sqrt x$ into
the single radical $\sqrt x$ over $\Q(\sqrt2)$. A negative content leaves
$\sqrt{-u}=\varepsilon i\sqrt u$ with $\varepsilon=\pm1$ locally constant. That
is the sign split of E6d over $\Q(i)$, still open (G-nf-11).

\part*{Part C. The gap log: what each extension closes}
\addcontentsline{toc}{part}{Part C}

Every gap the campaign closes is logged with non-trivial examples in
\code{tests/bronstein\_paper/examples.m}. Each example records:
\begin{itemize}
\item the ratchet site it removed (\code{tools/check\_bronstein\_completeness.py})
  and the section above that justifies it;
\item the release that closed it;
\item the verdict, and what the system did before.
\end{itemize}
The ctest \code{bronstein\_paper\_examples} checks every one. An antiderivative
must satisfy \code{PossibleZeroQ[D[F, x] - f]}, and a certificate must be
returned as one. \code{tools/bronstein\_paper\_table.py} regenerates the table
below, timings included (seconds, \code{Method -> "Bronstein"}, one core).

\begingroup\small\sloppy
% generated by tools/bronstein_paper_table.py from tests/bronstein_paper/examples.m -- do not edit
\begin{longtable}{>{\raggedright\arraybackslash}p{0.06\textwidth}>{\raggedright\arraybackslash}p{0.40\textwidth}>{\raggedright\arraybackslash}p{0.17\textwidth}>{\raggedright\arraybackslash}p{0.22\textwidth}r}
\toprule
Id & Input & Result (\S, version) & Before & s \\ \midrule
\endhead
P0-1 & \code{EIQ}: \code{Cos[E\textasciicircum{}x]}\newline{\footnotesize the lower monomial E\textasciicircum{}x reached the base-field solver in surface form; the empty result was read as a theorem} & \code{False} (--, v0.6545) & {\footnotesize False, by accident: the base solver could not form the equation} & 0.01 \\
P0-2 & \code{EIQ}: \code{E\textasciicircum{}x*Cos[E\textasciicircum{}x]}\newline{\footnotesize the same tower, elementary: Sin[E\textasciicircum{}x]} & \code{True} (--, v0.6545) & {\footnotesize True} & 0.00 \\
P1-1 & \code{Scope}: \code{Sqrt[Log[x] + Sqrt[Log[x]]]}\newline{\footnotesize nested radicals: outside one simple radical extension, reported as such} & implementation incomplete (--, v0.6545) & {\footnotesize unevaluated, silent} & 0.00 \\
P1-2 & \code{Scope}: \code{Sqrt[x] + Sqrt[x + 1]}\newline{\footnotesize Z/2 x Z/2: a genuine compositum} & implementation incomplete (--, v0.6545) & {\footnotesize unevaluated, silent} & 0.00 \\
P2-1 & \code{RDE}: \code{\{x\textasciicircum{}(-1), 2*(Log[x]/x\textasciicircum{}2)\}}\newline{\footnotesize the solution has degree dc - da + 1 in Log[x]: the bound an off-by-one first got wrong} & \code{Log[x]\textasciicircum{}2/x} (A2, v0.6546) & {\footnotesize undecided} & 0.00 \\
E6b-1 & \code{(Sqrt[x] + (x + 1)\textasciicircum{}(1/3))/(x*(x + 1))}\newline{\footnotesize y\textasciicircum{}6 = x\textasciicircum{}3 (x+1)\textasciicircum{}2, y = Sqrt[x] (x+1)\textasciicircum{}(1/3); answer in the user's radicals} & elementary (E6b, v0.6547) & {\footnotesize implementation incomplete} & 0.22 \\
E6b-2 & \code{Sqrt[x]*(x + 1)\textasciicircum{}(1/3)}\newline{\footnotesize Chebyshev's binomial case, certified by the Risch-Trager-Bronstein machinery on y\textasciicircum{}6 = x\textasciicircum{}3 (x+1)\textasciicircum{}2} & not elementary (certified) (E6b, v0.6547) & {\footnotesize implementation incomplete} & 0.02 \\
E6b-3 & \code{D[Sqrt[x - 1]*(x + 1)\textasciicircum{}(1/3)*E\textasciicircum{}x, x]}\newline{\footnotesize the compositum at the bottom of a transcendental tower (Shape B, E\textasciicircum{}x on top)} & elementary (E6b, v0.6547) & {\footnotesize implementation incomplete} & 0.17 \\
E6b-5 & \code{Sqrt[x]*(x + 1)\textasciicircum{}(1/3)*Log[x]}\newline{\footnotesize Log[x] over the compositum: certified by the coefficient Risch DE over the base} & not elementary (certified) (E6b, v0.6547) & {\footnotesize implementation incomplete} & 0.02 \\
E6b-6 & \code{(x*Sqrt[Log[x]]*(Log[x] + 1)\textasciicircum{}(1/3))\textasciicircum{}(-1)}\newline{\footnotesize Chebyshev's case in u = Log[x] (u\textasciicircum{}(-1/2) (1 + u)\textasciicircum{}(-1/3)), decided in the tower} & not elementary (certified) (E6b, v0.6547) & {\footnotesize implementation incomplete} & 0.03 \\
E6b-4 & \code{Sqrt[Log[x]]*((Log[x] + 1)\textasciicircum{}(1/3)/x)}\newline{\footnotesize the compositum over the transcendental tower Q(x)(Log[x])} & not elementary (certified) (E6b, v0.6547) & {\footnotesize implementation incomplete} & 0.03 \\
E6d-1 & \code{Sqrt[x\textasciicircum{}2*(x + 1)]/x\textasciicircum{}2 + (Sqrt[x + 1])\textasciicircum{}(-1)}\newline{\footnotesize Sqrt[x\textasciicircum{}2 (x+1)] = eps x Sqrt[x+1], eps = +-1 locally constant; the answer is a sum over sign regions} & elementary (E6d, v0.6549) & {\footnotesize implementation incomplete} & 0.08 \\
E6d-2 & \code{Sqrt[x\textasciicircum{}2*(x\textasciicircum{}3 + 1)] - x*Sqrt[x\textasciicircum{}3 + 1]}\newline{\footnotesize identically 0 where eps = 1, the elliptic -2 x Sqrt[x\textasciicircum{}3+1] where eps = -1} & not elementary (certified) (E6d, v0.6549) & {\footnotesize implementation incomplete} & 0.01 \\
E6d-3 & \code{Sqrt[x\textasciicircum{}2*(x\textasciicircum{}3 + 1)]*(Sqrt[x\textasciicircum{}3 + 1]/x)}\newline{\footnotesize the product is eps (x\textasciicircum{}3 + 1): rational on each region} & elementary (E6d, v0.6549) & {\footnotesize implementation incomplete} & 0.04 \\
E6d-4 & \code{Sqrt[Log[x]\textasciicircum{}2*(Log[x] + 1)] + Sqrt[Log[x] + 1]/x}\newline{\footnotesize over the tower: Log[x] Sqrt[Log[x] + 1] on the region eps = 1} & not elementary (certified) (E6d, v0.6549) & {\footnotesize implementation incomplete} & 0.02 \\
E6d-5 & \code{(1 + E\textasciicircum{}x\textasciicircum{}2)/2 + ((1 - E\textasciicircum{}x\textasciicircum{}2)/2)*Sqrt[x + 1]*Sqrt[(x + 1)\textasciicircum{}(-1)]}\newline{\footnotesize equals 1 on every open set; the sign -1 occurs only on the cut x < -1 and must never certify} & implementation incomplete (E6d, v0.6550) & {\footnotesize not elementary (v0.6549: false, a sign seen only on a branch cut)} & 0.01 \\
E6c-1 & \code{E\textasciicircum{}(x + Log[x]/2)/x}\newline{\footnotesize E\textasciicircum{}x/Sqrt[x] in disguise} & not elementary (certified) (E6c, v0.6547) & {\footnotesize undecided (a new hyperexponential monomial whose Risch DE needs a radical)} & 0.00 \\
E6c-4 & \code{E\textasciicircum{}(x\textasciicircum{}2 + Log[Log[x]]/2)}\newline{\footnotesize E\textasciicircum{}(x\textasciicircum{}2) Sqrt[Log[x]]: Shape A with the exponential below the radicand's monomial} & not elementary (certified) (E6c, v0.6547) & {\footnotesize undecided} & 0.01 \\
E6c-2 & \code{D[E\textasciicircum{}(x\textasciicircum{}2 + Log[x + 1]/3), x]}\newline{\footnotesize E\textasciicircum{}(x\textasciicircum{}2) (x+1)\textasciicircum{}(1/3): a radical at the bottom under an exponential} & elementary (E6c, v0.6547) & {\footnotesize undecided} & 0.01 \\
E6c-3 & \code{D[E\textasciicircum{}(Log[Log[x]]/2 + x), x]}\newline{\footnotesize Log[Log[x]] is rewritten away and must be pruned} & elementary (E6c, v0.6547) & {\footnotesize implementation incomplete (radical in the middle of the tower)} & 0.01 \\
E6a-1 & \code{Log[x\textasciicircum{}2]*(Sqrt[Log[x]]/x)}\newline{\footnotesize answer carries the branch difference Log[x\textasciicircum{}2] - 2 Log[x] explicitly} & elementary (E6a, v0.6547) & {\footnotesize implementation incomplete} & 0.01 \\
E6a-2 & \code{(Log[x\textasciicircum{}2] - 2*Log[x])*Sqrt[Log[x]]}\newline{\footnotesize zero on Re x > 0; certified by Hermite-Lindemann + specialisation (Thm E6a(3))} & not elementary (certified) (E6a, v0.6547) & {\footnotesize implementation incomplete} & 0.00 \\
E6a-3 & \code{Log[x*(x + 1)]*Sqrt[Log[x] + Log[x + 1]]*(x\textasciicircum{}(-1) + (x + 1)\textasciicircum{}(-1))}\newline{\footnotesize a product logarithm over two log monomials} & elementary (E6a, v0.6547) & {\footnotesize implementation incomplete} & 0.07 \\
E6a-4 & \code{(Log[E\textasciicircum{}x] - x)*Sqrt[Log[x]] + x\textasciicircum{}(-1)}\newline{\footnotesize the branch difference vanishes only on |Im x| < Pi; certified at a sample with large imaginary part} & not elementary (certified) (E6a, v0.6547) & {\footnotesize Log[x] from the Automatic cascade (wrong: PossibleZeroQ[Log[E\textasciicircum{}x] - x] is True)} & 0.01 \\
E3-1 & \code{(x*(1 + Log[x]\textasciicircum{}3)\textasciicircum{}(1/3))\textasciicircum{}(-1)}\newline{\footnotesize genus 1 (Fermat cubic over Q(x)(Log[x])); residues at theta = infinity} & elementary (E3, v0.6548) & {\footnotesize implementation incomplete} & 2.41 \\
E3-2 & \code{E\textasciicircum{}x/(1 + E\textasciicircum{}(3*x))\textasciicircum{}(1/3)}\newline{\footnotesize the exponential analogue: special places over theta = 0 and infinity} & elementary (E3, v0.6548) & {\footnotesize implementation incomplete} & 2.40 \\
E3-3 & \code{(x*(1 + Log[x]\textasciicircum{}4)\textasciicircum{}(1/4))\textasciicircum{}(-1)}\newline{\footnotesize m = 4, genus 3} & elementary (E3, v0.6548) & {\footnotesize implementation incomplete} & 5.68 \\
E3-4 & \code{D[(1 + Log[x] + Log[x]\textasciicircum{}3)\textasciicircum{}(1/3) + Log[Log[x] + (1 + Log[x] + Log[x]\textasciicircum{}3)\textasciicircum{}(1/3)], x]}\newline{\footnotesize algebraic part plus a log part with finite and special residues on y\textasciicircum{}3 = L\textasciicircum{}3 + L + 1} & elementary (E3, v0.6548) & {\footnotesize implementation incomplete} & 1.80 \\
E3-5 & \code{(1 + Log[x])/(x*(1 + Log[x] + Log[x]\textasciicircum{}3)\textasciicircum{}(1/3))}\newline{\footnotesize the residue divisor has infinite order (orders mod 13 and 37 disagree)} & not elementary (certified) (E3, v0.6548) & {\footnotesize implementation incomplete} & 1.18 \\
E3-6 & \code{D[x + Log[-E\textasciicircum{}x + (1 + E\textasciicircum{}(3*x))\textasciicircum{}(1/3)], x]}\newline{\footnotesize exponential tower, genus 1} & elementary (E3, v0.6548) & {\footnotesize implementation incomplete} & 4.93 \\
E41-1 & \code{Sqrt[x + I]/x}\newline{\footnotesize a curve over Q(i)(x); FriCAS agrees} & elementary (E4.1, v0.6551) & {\footnotesize implementation incomplete} & 0.07 \\
E41-2 & \code{D[Log[E\textasciicircum{}(I*x) + Sqrt[1 + E\textasciicircum{}(2*I*x)]], x]}\newline{\footnotesize Shape A over the exponential E\textasciicircum{}(I x): an inverse-hyperbolic of E\textasciicircum{}(I x)} & elementary (E4.1, v0.6551) & {\footnotesize implementation incomplete} & 0.09 \\
E41-3 & \code{(Sqrt[E\textasciicircum{}(2*I*x) + 1])\textasciicircum{}(-1)}\newline{\footnotesize residues at the special places of E\textasciicircum{}(I x) over Q(i)} & elementary (E4.1, v0.6551) & {\footnotesize implementation incomplete} & 0.09 \\
E41-4 & \code{(1 + I)*(Sqrt[x]/(x\textasciicircum{}2 + I))}\newline{\footnotesize residue field Q(i)(zeta\_8): i adjoined and written in the primitive element} & elementary (E4.1, v0.6551) & {\footnotesize implementation incomplete} & 0.09 \\
E41-5 & \code{Sqrt[Log[x] + I]}\newline{\footnotesize FriCAS: implementation incomplete (constant residues); Risch DE b' + b (1 + 1/(2 (u + I))) = 1 has no rational solution} & not elementary (certified) (E4.1, v0.6551) & {\footnotesize implementation incomplete} & 0.01 \\
\bottomrule
\end{longtable}

\endgroup

\paragraph{Open performance limits} (examples kept out of the table until they close):
\begin{itemize}
\item a three-radical compositum with Kummer group $\Z/30$,
  $\frac{d}{dx}\bigl(x^2\sqrt x\,(x+1)^{1/5}(x+2)^{1/3}\bigr)$: more than 300~s
  on the degree-30 curve;
\item $\frac{d}{dx}\bigl(\log(y+L)+2\log(y^2+Ly+L^2)\bigr)$ with $y^3=L^3+L+1$,
  $L=\log x$: 225~s.
\end{itemize}

\part*{Part D. Known open gaps}
\addcontentsline{toc}{part}{Part D}

Every known gap, in the campaign and in the lower layers its reviews uncovered,
is recorded in \code{tests/bronstein\_paper/open\_gaps.m}, with an example and its
current behaviour. The nightly ctest \code{bronstein\_open\_gaps\_stress} checks
that each entry still behaves as recorded. An entry that changes (a gap closed
or moved) fails the check until it is moved into the gap log (Part~C) together
with its fix. \code{tools/bronstein\_paper\_table.py --open-gaps} regenerates the
table below.

\begingroup\small\sloppy
% generated by tools/bronstein_paper_table.py --open-gaps from tests/bronstein_paper/open_gaps.m -- do not edit
\begin{longtable}{>{\raggedright\arraybackslash}p{0.08\textwidth}>{\raggedright\arraybackslash}p{0.42\textwidth}>{\raggedright\arraybackslash}p{0.13\textwidth}>{\raggedright\arraybackslash}p{0.27\textwidth}}
\toprule
Id & Gap and example & Now & Follow-up (found by) \\ \midrule
\endhead
G-perf-1 & a three-radical compositum (Kummer group Z/30): Trager on the degree-30 curve is too slow\newline\code{Sqrt[x]*(x + 1)\textasciicircum{}(1/5)*(x + 2)\textasciicircum{}(1/3)} & timeout & {\footnotesize profile Trager on high-degree Kummer curves; smaller presentations of the compositum (gap log (E6b), 2026-11-14)} \\
G-perf-2 & two log terms with special residues on a genus-1 cubic: 225 s\newline\code{Together[D[Log[Log[x] + (1 + Log[x] + Log[x]\textasciicircum{}3)\textasciicircum{}(1/3)] + 2*Log[Log[x]\textasciicircum{}2 + Log[x]*(1 + Log[x] + Log[x]\textasciicircum{}3)\textasciicircum{}(1/3) + (1 + Log[x] + Log[x]\textasciicircum{}3)\textasciicircum{}(2/3)], x]]} & timeout & {\footnotesize profile the E3 torsion / principal-generator steps (E3 tests, 2026-11-14)} \\
G-perf-3 & the fourth root of the tangent times out (13 of 120 random m = 3, 4 components also time out at 30 s)\newline\code{Tan[x]\textasciicircum{}(1/4)} & timeout & {\footnotesize profile m = 4 over Q(i); the random-component sweep of E6e (review of E4.2; E6e sweep)} \\
G-sign-1 & radicals equal up to a sign seen only on a branch cut: undecided (E6d)\newline\code{(1 + E\textasciicircum{}x\textasciicircum{}2)/2 + ((1 - E\textasciicircum{}x\textasciicircum{}2)/2)*Sqrt[x + 1]*Sqrt[(x + 1)\textasciicircum{}(-1)]} & incomplete & {\footnotesize prove eps = 1 on every open set (no open region with eps = -1) and integrate f\_+ alone (review of E6d)} \\
G-sign-2 & Sqrt[Tan] + Sqrt[Cot]: the sign split's own check fails (an Internal error)\newline\code{Sqrt[Tan[x]] + Sqrt[Cot[x]]} & internal & {\footnotesize debug SignSplitCheckFailed over Q(i) (review of E4.2)} \\
G-trig-1 & TrigToExp writes ArcCosh with Sqrt[x - 1] Sqrt[x + 1]: two radicals (E6d over a constant-sign split)\newline\code{ArcCosh[x]/Sqrt[x\textasciicircum{}2 - 1]} & incomplete & {\footnotesize a front-end rewrite Sqrt[x - 1] Sqrt[x + 1] = eps Sqrt[x\textasciicircum{}2 - 1] (E6d) (review of E4.2)} \\
G-qi-1 & residue at infinity of a logarithm: the integral in k has a logarithm with non-rational coefficient (Q(i))\newline\code{Sqrt[I + Log[x]\textasciicircum{}2]} & incomplete & {\footnotesize brLimitedIntegrate over Q(i) with algebraic log coefficients (Phase 3 LimitedIntegrate) (Phase 4.1b battery)} \\
G-nf-1 & number-field constants (Sqrt[2], ...): Tower:algebraic-constant\newline\code{Sqrt[2]*Sqrt[Log[x]]*Log[x + 1]} & incomplete & {\footnotesize Phase 4.3 (plan)} \\
G-nf-2 & radicals equal over a larger constant field (Sqrt[2 x] = Sqrt[2] Sqrt[x])\newline\code{Sqrt[2*x] + Sqrt[x]} & incomplete & {\footnotesize Phase 4.3 (E6b/E6d over a number field) (E6b)} \\
G-red-1 & two reducible logarithms: kappa values may be algebraically dependent (E6a needs one)\newline\code{(Log[x\textasciicircum{}2] - 2*Log[x])*(Log[x\textasciicircum{}3] - 3*Log[x])*Sqrt[Log[x]]} & incomplete & {\footnotesize E6a for several kappa (their values are logs of algebraic numbers: Baker's theorem on linear forms) (E6a)} \\
G-mid-1 & a radical in the middle of the tower\newline\code{Log[1 + Sqrt[Log[x]]]/x} & incomplete & {\footnotesize Phase 5 (E1, E2) (plan)} \\
G-param-1 & symbolic parameters\newline\code{Sqrt[a + Log[x]]/x} & incomplete & {\footnotesize Phase 4.4 (E4: generic verdicts) (plan)} \\
G-pres-1 & answers over Q(i) carry Root objects of the residue field (no radicals, no real form)\newline\code{Sqrt[Tan[x]]} & presentation & {\footnotesize ToRadicals on residue-field generators; a verified real form for Q(i) answers of real integrands (Phase 4.2)} \\
G-core-1 & PossibleZeroQ false zero: Log[E\textasciicircum{}x] - x is -2 Pi I for Im x in (Pi, 3 Pi]\newline\code{PossibleZeroQ[Log[E\textasciicircum{}x] - x]} & True & {\footnotesize PossibleZeroQ: sample off the strip |Im x| < Pi for Log of exponentials (review of E6a)} \\
G-core-2 & the Automatic cascade returns Indeterminate (Bronstein certifies it)\newline\code{Quiet[Integrate[(Log[x\textasciicircum{}2] - 2*Log[x])*Exp[x\textasciicircum{}2], x]]} & Indeterminate & {\footnotesize find the stage that returns Indeterminate (review of E6a)} \\
G-core-3 & an uncaught Throw[..., "PIM"] from the Automatic cascade\newline\code{Head[Quiet[Integrate[Exp[x + Log[x\textasciicircum{}2]/4], x]]]} & Throw & {\footnotesize catch the PIM throw at its stage boundary (review of E6c)} \\
G-core-4 & Log[x\textasciicircum{}2] is not its own monomial in the C transcendental integrator (the retry fixed in v0.6547 no longer answers wrongly)\newline\code{Head[Integrate`RischTranscendental[Log[x\textasciicircum{}2]/x, x]]} & Integrate`RischTranscendental & {\footnotesize tower builder: a logarithm of a power as a monomial, or the kappa split of E6a in C (Phase 3 part 1)} \\
G-core-5 & machine ArcSech at negative reals disagrees with its exact form (latent: no integrand reaches it)\newline\code{Sign[Im[ArcSech[-3.7]]] == Sign[Im[N[ArcSech[-37/10]]]]} & False & {\footnotesize fix the machine ArcSech branch for x < 0 (review of E4.2)} \\
\midrule
\multicolumn{4}{l}{\emph{Gaps without a single example}} \\
ratchet & \multicolumn{2}{>{\raggedright\arraybackslash}p{0.55\textwidth}}{41 throwUnimplemented sites and 12 uncertain C Risch DE paths (tools/check\_bronstein\_completeness.py, BASELINE\_M / BASELINE\_C)} & {\footnotesize Phases 2-5} \\
Phase 2 & \multicolumn{2}{>{\raggedright\arraybackslash}p{0.55\textwidth}}{ParamRischDE and LimitedIntegrate with new logarithms (Ch. 7); rewrite dependent Log/Exp towers; rt\_field\_integrate completeness (m > 1 new logs, mixed resultants)} & {\footnotesize Phase 2 remainder} \\
E6e & \multicolumn{2}{>{\raggedright\arraybackslash}p{0.55\textwidth}}{LogPart:weights at an exponential's special places (Bronstein's Def 2.21 residues): the weight theorem not yet written for them} & {\footnotesize Phase 3} \\
E5 & \multicolumn{2}{>{\raggedright\arraybackslash}p{0.55\textwidth}}{several ramified places above one special point (Divisors:multi-ramified)} & {\footnotesize Phase 3} \\
4.1 & \multicolumn{2}{>{\raggedright\arraybackslash}p{0.55\textwidth}}{Trager:residue-field-QI: the tower method of residueField over Q(i)} & {\footnotesize Phase 4.1} \\
certificates & \multicolumn{2}{>{\raggedright\arraybackslash}p{0.55\textwidth}}{certificates without an external cross-check (FriCAS times out or reports incomplete): Sqrt[Log[x] + I], 1/((x - I) Sqrt[x\textasciicircum{}3 + 1]), (1 + Log[x])/(x (1 + Log[x] + Log[x]\textasciicircum{}3)\textasciicircum{}(1/3))} & {\footnotesize an independent checker (another CAS or a proof certificate)} \\
tests & \multicolumn{2}{>{\raggedright\arraybackslash}p{0.55\textwidth}}{integrate\_cascade\_trager\_bronstein\_tests:89 emits a message only under full parallel ctest load; repl\_pty\_test and image\_tests hit their alarms under load} & {\footnotesize test infrastructure} \\
design & \multicolumn{2}{>{\raggedright\arraybackslash}p{0.55\textwidth}}{goodPrimes searches primes without a bound: data that no prime is good for (inconsistent input) would loop; characteristic-0 prechecks guard the known cases} & {\footnotesize a proof-backed bound or an Internal error} \\
\bottomrule
\end{longtable}

\endgroup

\end{document}
